% ID: SEL-D0022-OB-E02
% ESTADO: cotejada. Segunda pasada visual completa y borrador compilado revisados.
% Fuente: Selectividad/matematicas_aplicadas/reserva_a_EXAMEN 5 -ANDALUCÍA.pdf
% SHA-256: 668689642ca997cf817c9f9f6608d2835efe898e3162b050d7b22b51f225d363
% Página: 2; opción B; ejercicio original 2.

Los beneficios de una empresa, en miles de euros, han evolucionado en los 25 años de su existencia según una función del tiempo, en años, dada por la siguiente expresión:
\[
B(t)=\begin{cases}
4t & \text{si }0\leq t<10\\[4pt]
-\dfrac{1}{5}t^2+8t-20 & \text{si }10\leq t\leq 25
\end{cases}
\]
\begin{enumerate}[label=\alph*)]
\item \textbf{(1 punto)} Estudie la continuidad y derivabilidad de $B$ en el intervalo $[0,25]$.
\item \textbf{(1 punto)} Estudie la monotonía de esta función y determine en qué año fueron mayores los beneficios de esta empresa y cuál fue su beneficio máximo.
\item \textbf{(0.5 puntos)} Represente gráficamente esta función.
\end{enumerate}
